\chapter{Metody analizy statystycznej strumienia bitów losowych}

Rozdział łączy omówienie miar entropii, pojęcia losowości oraz metod analizy
statystycznej. Połączenie to jest celowe: entropii nie da się omówić bez
wskazania konkretnej miary, a każda praktycznie stosowana miara entropii jest
estymowana z ciągu bitów procedurami statystycznymi. Metody analizy
statystycznej nie są zatem oddzielnym zagadnieniem, lecz aparatem pomiarowym
entropii.

\section{Miary entropii}

Niech $X$ będzie dyskretną zmienną losową o rozkładzie $p_1,\dots,p_n$.
Stosowane w ocenie generatorów miary tworzą rodzinę entropii Rényiego
rzędu $\alpha$ \cite{renyi1961}:
%
\begin{equation}
  H_\alpha(X) = \frac{1}{1-\alpha}\log_2\sum_{i=1}^{n} p_i^{\alpha},
  \qquad \alpha > 0,\ \alpha \neq 1 .
  \label{eq:renyi}
\end{equation}
%
Szczególne przypadki wyrażenia \eqref{eq:renyi} mają w kontekście generatorów
odrębną interpretację:

\begin{itemize}
  \item \textbf{entropia Shannona} ($\alpha \to 1$):
        \begin{equation}
          H_1(X) = -\sum_{i=1}^{n} p_i \log_2 p_i ,
          \label{eq:shannon}
        \end{equation}
        średnia liczba bitów informacji na symbol, miara przyjęta
        w AIS-31 \cite{killmann2011ais31};
  \item \textbf{entropia kolizji} ($\alpha = 2$):
        \begin{equation}
          H_2(X) = -\log_2\sum_{i=1}^{n} p_i^2 ,
          \label{eq:kolizji}
        \end{equation}
        związana z prawdopodobieństwem kolizji, wykorzystywana
        w estymatorach i w twierdzeniach o ekstrakcji losowości;
  \item \textbf{min-entropia} ($\alpha \to \infty$):
        \begin{equation}
          H_\infty(X) = -\log_2 \max_{i} p_i ,
          \label{eq:minentropia}
        \end{equation}
        miara najgorszego przypadku, odpowiadająca prawdopodobieństwu
        odgadnięcia wartości przy jednej próbie, podstawowa miara
        w SP 800-90B \cite{turan2018sp80090b}, normie określającej wymagania stawiane źródłom entropii stosowanym w generatorach bitów losowych oraz procedury ich walidacji.
\end{itemize}

Zachodzi porządek
%
\begin{equation}
  H_\infty(X) \le H_2(X) \le H_1(X) \le \log_2 n ,
  \label{eq:porzadek}
\end{equation}
%
przy czym w zastosowaniach kryptograficznych istotna jest właśnie miara
najsłabsza. Ilustruje to przykład źródła Bernoulliego o obciążeniu $p$,
zestawiony w tabeli \ref{tab:entropia-obciazenie}.

\begin{table}[htbp]
  \centering
  \caption{Entropia Shannona i min-entropia źródła Bernoulliego
           w funkcji obciążenia}
  \label{tab:entropia-obciazenie}
  \begin{tabular}{ccc}
    \hline
    $p = P(X = 1)$ & $H_1$ [bit/bit] & $H_\infty$ [bit/bit] \\
    \hline
    $0{,}50$ & $1{,}000$ & $1{,}000$ \\
    $0{,}55$ & $0{,}993$ & $0{,}862$ \\
    $0{,}60$ & $0{,}971$ & $0{,}737$ \\
    $0{,}70$ & $0{,}881$ & $0{,}515$ \\
    $0{,}80$ & $0{,}722$ & $0{,}322$ \\
    \hline
  \end{tabular}
\end{table}

Dla małego obciążenia $p = \tfrac{1}{2} + \varepsilon$ zachodzą przybliżenia
%
\begin{equation}
  H_1 \approx 1 - \frac{2}{\ln 2}\,\varepsilon^2 ,
  \qquad
  H_\infty \approx 1 - \frac{2}{\ln 2}\,\varepsilon .
  \label{eq:przyblizenia}
\end{equation}
%
Entropia Shannona maleje więc kwadratowo, a min-entropia liniowo. Pierwsza
z miar jest zatem nadmiernie \emph{optymistyczna}: źródło o obciążeniu
$p = 0{,}6$ wygląda w niej niemal idealnie ($0{,}971$), a w kategoriach
odgadywalności traci ponad jedną czwartą siły. Dla klucza 128-bitowego
złożonego z takich bitów efektywna siła wynosi około 94 bity.

\subsection{Entropia źródeł z pamięcią}

W przypadku źródeł z pamięcią rozkład brzegowy nie wystarcza - konieczne
jest posługiwanie się \textbf{entropią warunkową}
$H(X_n \mid X_{n-1},\dots,X_1)$ oraz \textbf{stopą entropii} procesu
stacjonarnego \cite{cover2006elements}:
%
\begin{equation}
  \bar{H} = \lim_{n \to \infty} \frac{1}{n} H(X_1,\dots,X_n),
  \label{eq:stopa-entropii}
\end{equation}
%
a w praktyce modelami Markowa skończonego rzędu. Zaniedbanie zależności
czasowych jest jednym z najczęstszych błędów oceny generatorów sprzętowych:
ciąg o idealnym rozkładzie brzegowym i silnej autokorelacji ma stopę entropii
znacznie mniejszą od jedności.

Znaczenie tego rozróżnienia dla praktyki pomiarowej jest bezpośrednie.
Rozkład brzegowy bada się miarą obciążenia (bias) i testem monobitowym,
natomiast obecność pamięci - funkcją autokorelacji oraz rozkładem słów
wielobitowych. Są to dwie \emph{niezależne} rodziny własności: strumień może
przejść pomiarowo pierwszą z nich i całkowicie zawieść w drugiej. Dlatego
w części doświadczalnej niniejszej pracy oba aspekty analizowane są osobno,
a nie jako pojedyncza ocena jakości strumienia.

\section{Pojęcie ciągu losowego}

Miary wprowadzone powyżej opisują \emph{źródło}. Osobnym pytaniem jest, co
znaczy, że konkretny \emph{ciąg} jest losowy. Historycznie ukształtowały się
trzy podejścia do formalizacji tego pojęcia.

\paragraph{Podejście częstościowe (statystyczne).}
Zapoczątkowane przez Richarda von Misesa \cite{vonmises1919} pojęcie
\emph{Kollektiv} (ang. \emph{collective}) wymaga, aby częstości symboli
w ciągu i we wszystkich jego podciągach wybranych dopuszczalną regułą były
zbieżne do wartości oczekiwanych. Zbiór dopuszczalnych reguł wyboru
doprecyzował Church \cite{church1940}, ograniczając go do reguł obliczalnych.
Jean Ville \cite{ville1939} wykazał, że warunki tego typu nie wystarczają:
istnieją ciągi spełniające je, a jednocześnie wykazujące jawną regularność
(na przykład sumy częściowe zawsze po jednej stronie wartości oczekiwanej).
Podejście to jest natomiast bezpośrednim źródłem konstrukcji testów
statystycznych, ponieważ każdy test bada jedną wybraną własność częstościową.

\paragraph{Podejście złożonościowe.}
Według Kołmogorowa \cite{kolmogorov1965three} i Martina-Löfa
\cite{martinlof1966} ciąg jest losowy, jeśli nie istnieje jego opis krótszy
od samego ciągu, co oznacza, że jego złożoność algorytmiczna jest bliska
jego długości. Złożoność Kołmogorowa jest nieobliczalna, więc nie istnieje
uniwersalny, skończony test losowości. Ma to fundamentalną konsekwencję
praktyczną: badanie generatora może jedynie falsyfikować hipotezę
o losowości, nigdy natomiast jej nie potwierdzi.

\paragraph{Podejście obliczeniowe (nieprzewidywalność).}
W kryptografii przyjmuje się kryterium nieodróżnialności obliczeniowej od
ciągu idealnego \cite{yao1982theory}, które jest równoważne testowi
następnego bitu \cite{blum1984how,goldreich2001foundations}. Generator jest
bezpieczny, jeśli żaden algorytm o ograniczonych zasobach nie przewiduje
kolejnego bitu z prawdopodobieństwem istotnie większym niż $1/2$. Do tego
dochodzą wymagania dotyczące kompromitacji stanu wewnętrznego. Odporność na
cofanie (ang. \emph{backtracking resistance}) oznacza, że ujawnienie stanu
nie pozwala odtworzyć bitów wcześniejszych. Odporność na przewidywanie
(ang. \emph{prediction resistance}) oznacza, że ujawnienie stanu nie pozwala
przewidzieć bitów przyszłych po dostarczeniu do generatora świeżej entropii
\cite{barker2015sp80090a}. Podobne wymagania, nazywane tajnością wprzód
i wstecz (ang. \emph{forward/backward secrecy}), formułuje AIS~20
\cite{killmann2011ais31}.

\subsection{Ograniczenia testów statystycznych a ocena źródła}

Warto na tym tle wskazać ograniczenie testów statystycznych. Rozważmy rejestr
przesuwny ze sprzężeniem liniowym (ang. \emph{Linear Feedback Shift
Register}, LFSR). Jest to układ złożony z $m$ komórek pamięci, w którym
w każdym takcie zawartość rejestru przesuwa się o jedną pozycję, a nowy bit
wprowadzany na jego początek jest sumą modulo~2 (operacją XOR) wybranych
komórek, zwanych odczepami. Generowany ciąg $s_t$ spełnia zależność
rekurencyjną
%
\begin{equation}
  s_{t+m} = \bigoplus_{j=0}^{m-1} c_j\, s_{t+j},
  \qquad c_j \in \{0,1\},
  \label{eq:lfsr}
\end{equation}
%
w której współczynniki $c_j$ wskazują położenie odczepów i wyznaczają
wielomian charakterystyczny $P(x) = x^m + \sum_{j=0}^{m-1} c_j x^j$ nad ciałem
$\mathrm{GF}(2)$. Jeżeli $P(x)$ jest wielomianem pierwotnym, to rejestr
zainicjowany dowolnym stanem niezerowym generuje ciąg o maksymalnym okresie
$2^m - 1$, zwany ciągiem $m$ \cite{golomb1982shift}. Ciąg taki ma własności,
których oczekuje się od ciągu losowego. W jednym okresie liczba jedynek
różni się od liczby zer o jeden, rozkład długości serii jest zgodny
z rozkładem idealnym, a autokorelacja przyjmuje dwie wartości, zbliżone do
zera dla wszystkich przesunięć niezerowych. Ciąg przechodzi więc większość
klasycznych testów statystycznych.

Jednocześnie jest on całkowicie deterministyczny. Cały ciąg wynika
z $m$-bitowego stanu początkowego i znanych odczepów, a algorytm
Berlekampa-Masseya \cite{massey1969shift} odtwarza strukturę rejestru już
na podstawie $2m$ kolejnych bitów wyjścia. Entropia takiego ciągu nie
przekracza $m$ bitów, niezależnie od jego długości, ponieważ każdy bit
wynika z pierwotnego stanu. Na przykład ciąg o długości $10^6$ bitów
wytworzony przez rejestr 32-bitowy zawiera co najwyżej 32 bity entropii,
a entropia przypadająca na bit maleje do zera wraz z długością ciągu.

Wniosek ma znaczenie dla dalszej części pracy. Test statystyczny ocenia
\emph{ciąg}, natomiast entropia jest własnością \emph{mechanizmu}, który
ten ciąg wytwarza. Wynik pozytywny nie dowodzi więc nieprzewidywalności,
a ocena źródła musi opierać się na modelu jego działania i na analizie
surowych próbek przed jakąkolwiek obróbką deterministyczną
\cite{turan2018sp80090b}.

\section{Aparat testowania hipotez}

Każdy test statystyczny losowości ma tę samą strukturę formalną. Formułuje
się hipotezę zerową
%
\begin{equation}
  H_0: \text{badany ciąg jest realizacją źródła i.i.d. o } P(0) = P(1) = \tfrac{1}{2},
  \label{eq:h0}
\end{equation}
%
wyznacza statystykę testową $T$, której rozkład przy prawdziwości $H_0$ jest
znany (najczęściej asymptotycznie normalny lub $\chi^2$), i oblicza się
wartość $p$ jako prawdopodobieństwo uzyskania wyniku co najmniej tak
skrajnego jak zaobserwowany:
%
\begin{equation}
  p = P\big(|T| \ge |t_{\mathrm{obs}}| \;\big|\; H_0\big).
  \label{eq:pvalue}
\end{equation}
%
Hipotezę odrzuca się, gdy $p < \alpha$, gdzie $\alpha$ jest przyjętym
poziomem istotności (w pakiecie SP 800-22 standardowo $\alpha = 0{,}01$).
Interpretacja wyniku wymaga trzech zastrzeżeń, istotnych dla poprawnego
odczytania rezultatów pomiarowych.

Po pierwsze, $p \ge \alpha$ oznacza wyłącznie brak podstaw do odrzucenia
$H_0$, nigdy jej potwierdzenia - co jest bezpośrednią konsekwencją
nieobliczalności złożoności Kołmogorowa omówionej wcześniej.

Po drugie, przy wielokrotnym testowaniu odrzucenia pojawiają się z
częstością $\alpha$ nawet dla źródła idealnego. Dla $k$ niezależnych testów
oczekiwana liczba fałszywych odrzuceń wynosi $k\alpha$, a przy $\alpha = 0{,}01$
i $k = 100$ podciągach jedno odrzucenie jest wynikiem \emph{typowym}, nie
sygnałem błędu. Poprawnym kryterium jest więc nie „brak odrzuceń”, lecz
zgodność \emph{rozkładu} uzyskanych wartości $p$ z rozkładem jednostajnym
$U(0,1)$, badana testem $\chi^2$ lub testem Kołmogorowa--Smirnowa.

Po trzecie, moc testu rośnie z długością ciągu. Odstępstwo o ustalonej
wielkości - np. obciążenie $\varepsilon$ - pozostaje niewykrywalne dla
próby zbyt małej i jest wykrywane niemal zawsze dla próby dostatecznie
dużej. Wielkość odstępstwa (rozmiar efektu) i jego istotność statystyczna
są zatem pojęciami rozłącznymi: przy strumieniu rzędu $10^8$ bitów wykrywalne
stają się obciążenia praktycznie nieistotne z punktu widzenia entropii,
i odwrotnie --- krótki rekord nie pozwala orzec o braku obciążenia.

\section{Analiza rozkładu brzegowego: obciążenie i test monobitowy}

Najprostszą badaną własnością jest zgodność częstości jedynek z wartością
oczekiwaną $1/2$. Dla ciągu $b_1,\dots,b_N$ definiuje się częstość empiryczną
i obciążenie
%
\begin{equation}
  \hat{p} = \frac{1}{N}\sum_{i=1}^{N} b_i ,
  \qquad
  \varepsilon = \hat{p} - \tfrac{1}{2}.
  \label{eq:bias}
\end{equation}
%
W teście monobitowym (test częstości, \emph{frequency test}) stosuje się
statystykę opartą na sumie ciągu przekształconego do postaci $\pm 1$,
$S_N = \sum_{i=1}^{N}(2b_i - 1)$, dla której przy prawdziwości $H_0$
z centralnego twierdzenia granicznego wynika
%
\begin{equation}
  Z = \frac{|S_N|}{\sqrt{N}} \;\xrightarrow{\;d\;}\; |\mathcal{N}(0,1)| ,
  \qquad
  p = \operatorname{erfc}\!\left(\frac{|S_N|}{\sqrt{2N}}\right).
  \label{eq:monobit}
\end{equation}
%
Odchylenie standardowe estymatora $\hat{p}$ wynosi $1/(2\sqrt{N})$, co
wyznacza rozdzielczość pomiaru: przy $N = 10^6$ bitów niepewność oszacowania
obciążenia jest rzędu $5\cdot 10^{-4}$, a obciążenia mniejsze od tej wartości
nie są odróżnialne od fluktuacji statystycznej. Ma to znaczenie przy
porównywaniu obciążeń poszczególnych kanałów generatora - różnica między
nimi jest istotna tylko wtedy, gdy przekracza niepewność wynikającą z
długości rekordu.

Powiązanie z miarami entropii jest bezpośrednie: podstawienie $\hat{p}$ do
wzorów \eqref{eq:shannon} i \eqref{eq:minentropia} daje estymatory entropii
Shannona i min-entropii \emph{przy założeniu niezależności bitów}.
Oszacowania te stanowią górne ograniczenie faktycznej entropii strumienia
- każda pamięć źródła może je tylko obniżyć - i dlatego wymagają
uzupełnienia analizą korelacji.

\section{Analiza zależności czasowych: autokorelacja}

Do wykrywania pamięci źródła służy funkcja autokorelacji. Dla ciągu bitów
o wartościach $\{0,1\}$ i średniej $\hat{p}$ definiuje się znormalizowany
współczynnik autokorelacji dla opóźnienia (przesunięcia) $\tau$:
%
\begin{equation}
  r(\tau) = \frac{\displaystyle\sum_{i=1}^{N-\tau}(b_i - \hat{p})(b_{i+\tau} - \hat{p})}
                 {\displaystyle\sum_{i=1}^{N}(b_i - \hat{p})^2},
  \qquad \tau = 1,2,\dots,\tau_{\max}.
  \label{eq:autokorelacja}
\end{equation}
%
Dla źródła i.i.d. i dużego $N$ estymator $r(\tau)$ ma rozkład w przybliżeniu
normalny o zerowej średniej i odchyleniu standardowym $1/\sqrt{N}$, skąd
przedział ufności na poziomie $1-\alpha$ ma postać
%
\begin{equation}
  |r(\tau)| \le \frac{z_{1-\alpha/2}}{\sqrt{N}} ,
  \label{eq:ci-autokorelacja}
\end{equation}
%
co dla $\alpha = 0{,}05$ daje granicę $\approx 1{,}96/\sqrt{N}$. Wartości
przekraczające ten próg wskazują na obecność zależności o danym opóźnieniu.

Interpretacja przebiegu $r(\tau)$ pozwala rozróżnić kilka typów defektu
źródła fizycznego:

\begin{itemize}
  \item \textbf{pojedynczy pik przy $\tau = 1$} - pamięć krótkozasięgowa,
        typowa dla zbyt szybkiego próbkowania względem czasu relaksacji
        źródła; kolejne próbki nie są od siebie niezależne;
  \item \textbf{składowa periodyczna} $r(\tau) \sim \cos(2\pi\tau/T)$ -
        wskazuje na sprzężenie z sygnałem okresowym (zegar, zasilanie,
        zakłócenie zewnętrzne) i jest najpoważniejszym defektem, gdyż
        oznacza deterministyczny wkład do strumienia;
  \item \textbf{wykładniczy zanik} $r(\tau) \sim e^{-\tau/\tau_c}$ -
        pamięć o skończonym czasie korelacji $\tau_c$; sugeruje decymację
        strumienia z krokiem większym niż $\tau_c$;
  \item \textbf{brak struktury w granicach} \eqref{eq:ci-autokorelacja} -
        brak podstaw do odrzucenia niezależności dla badanego zakresu
        opóźnień.
\end{itemize}

Należy podkreślić, że zerowa autokorelacja nie implikuje niezależności:
$r(\tau)$ mierzy wyłącznie zależność liniową, a wykrycie zależności
wyższego rzędu wymaga analizy rozkładu słów wielobitowych. Przy badaniu
$\tau_{\max}$ opóźnień jednocześnie obowiązuje przy tym korekta na
wielokrotne testowanie omówiona w podrozdziale~\ref{sec:aparat-hipotez}.

\section{Analiza słów wielobitowych i mapy dwuwymiarowe}

Uogólnieniem analizy pojedynczych bitów jest badanie rozkładu słów
$k$-bitowych. Grupując strumień w słowa $w_j \in \{0,\dots,2^k-1\}$ i
zestawiając liczby ich wystąpień $O_m$ z wartością oczekiwaną
$E = M/2^k$ (gdzie $M$ jest liczbą słów), otrzymuje się statystykę
%
\begin{equation}
  \chi^2 = \sum_{m=0}^{2^k-1} \frac{(O_m - E)^2}{E}
  \;\sim\; \chi^2_{2^k-1},
  \label{eq:chi2-slowa}
\end{equation}
%
która pozwala testować jednostajność rozkładu słów (test serii, w wariantach
literaturowych określany też jako \emph{serial test} lub \emph{poker test}).
Bezpośrednie oszacowanie entropii z częstości słów, $\hat{H} = -\sum_m
\hat{q}_m \log_2 \hat{q}_m$, jest obciążone w dół i wymaga korekty
Millera--Madowa rzędu $(2^k-1)/(2M\ln 2)$; z tego względu wiarygodna
estymacja wymaga $M \gg 2^k$, co ogranicza praktycznie osiągalne $k$.

Uzupełnieniem analizy liczbowej jest reprezentacja graficzna w postaci
\textbf{mapy dwuwymiarowej}, konstruowanej przez naniesienie na płaszczyznę
par kolejnych słów
%
\begin{equation}
  \big(w_j,\, w_{j+d}\big), \qquad j = 1,\dots,M-d,
  \label{eq:mapa2d}
\end{equation}
%
gdzie $d$ jest przyjętym opóźnieniem. Metoda ta jest odpowiednikiem
rekonstrukcji przestrzeni fazowej metodą opóźnień czasowych, stosowanej
w analizie układów dynamicznych: dla źródła idealnego punkty wypełniają
kwadrat $[0,2^k)^2$ równomiernie, a dla źródła zawierającego składową
deterministyczną - układają się w struktury geometryczne. Wartość
diagnostyczna map wynika stąd, że ujawniają one zależności
\emph{nieliniowe}, wobec których funkcja autokorelacji pozostaje ślepa,
a ich typowe wzorce mają bezpośrednią interpretację układową:

\begin{itemize}
  \item \textbf{linie ukośne} - zależność funkcyjna między słowami
        odległymi o $d$, czyli deterministyczne odwzorowanie stanu;
  \item \textbf{prążki poziome lub pionowe} - niejednorodność rozkładu
        brzegowego słów (uprzywilejowane wartości), wykrywalna również
        testem \eqref{eq:chi2-slowa};
  \item \textbf{zagęszczenie wokół przekątnej} - dodatnia korelacja
        kolejnych słów, odpowiadająca pamięci źródła;
  \item \textbf{regularna krata} - okresowość strumienia, wskazująca na
        skończony i krótki cykl efektywny.
\end{itemize}

Ocena wzrokowa mapy stanowi wstępny etap diagnostyki; sformalizowaniem tej
samej informacji jest omawiany wyżej test $\chi^2$ dla par słów, o liczbie
stopni swobody $2^{2k}-1$.

\section{Kondycjonowanie strumienia i metoda von Neumanna}

Ponieważ surowy strumień ze źródła fizycznego jest z reguły obciążony,
stosuje się etap kondycjonowania (\emph{wybielania}), którego zadaniem jest
usunięcie odstępstw od rozkładu jednostajnego \emph{bez} zwiększania
faktycznej ilości niesionej entropii. Klasycznym algorytmem tego typu jest
korektor von Neumanna \cite{vonneumann1951}: bity grupuje się w rozłączne
pary $(b_{2i}, b_{2i+1})$ i stosuje odwzorowanie
%
\begin{equation}
  (0,1) \mapsto 0, \qquad (1,0) \mapsto 1, \qquad
  (0,0),\,(1,1) \mapsto \text{odrzucenie pary}.
  \label{eq:von-neumann}
\end{equation}
%
Dla źródła Bernoulliego o obciążeniu $p$ i niezależnych bitach wynik jest
ciągiem idealnie symetrycznym niezależnie od wartości $p$, co wynika z
równości $P(0,1) = P(1,0) = p(1-p)$. Prawdopodobieństwo zachowania pary
wynosi $2p(1-p)$, więc wydajność metody, mierzona liczbą bitów wyjściowych
na bit wejściowy, jest równa
%
\begin{equation}
  \eta(p) = p(1-p) \;\le\; \tfrac{1}{4},
  \label{eq:wydajnosc-vn}
\end{equation}
%
z maksimum $0{,}25$ osiąganym dla $p = 0{,}5$. Nawet dla źródła idealnego
odrzucane są zatem trzy czwarte bitów, a wydajność maleje wraz ze wzrostem
obciążenia. Warto porównać tę wartość z granicą teoretyczną: entropia
źródła wynosi $H_1(p)$ bitów na bit, więc metoda von Neumanna wykorzystuje
jedynie ułamek $\eta(p)/H_1(p)$ dostępnej entropii (dla $p = 0{,}5$ -
25\%), podczas gdy ekstraktory optymalne zbliżają się do granicy $H_1(p)$.

Ograniczenia metody są istotne dla interpretacji wyników pomiarowych.
Poprawność dowodu opiera się na założeniu niezależności bitów w parze:
przy obecności autokorelacji symetria $P(0,1) = P(1,0)$ przestaje zachodzić
i wybielanie nie usuwa obciążenia w pełni. Metoda nie usuwa też korelacji
- eliminuje odstępstwo rozkładu brzegowego, natomiast struktura zależności
czasowych może zostać zachowana, a przez nieregularne odrzucanie par nawet
przekształcona w sposób trudny do przewidzenia. Z tego względu porównanie
autokorelacji i rozkładu słów \emph{przed} i \emph{po} wybielaniu jest
osobnym elementem oceny generatora, a nie formalnością: poprawa obciążenia
przy niezmienionej strukturze korelacyjnej wskazuje, że wybielanie
maskuje objaw, nie usuwając przyczyny.

W zastosowaniach wymagających odporności na korelacje długozasięgowe
stosuje się w zamian ekstraktory kryptograficzne (uniwersalne funkcje
haszujące, konstrukcje zgodne z SP 800-90C), które przy poprawnie
oszacowanej min-entropii wejścia gwarantują nieodróżnialność wyjścia od
ciągu idealnego z zadanym zapasem bezpieczeństwa.

\section{Standaryzowane pakiety testów}

Opisane wyżej testy stanowią rdzeń analizy własnej, prowadzonej z pełną
kontrolą nad parametrami pomiaru. Ich uzupełnieniem są standaryzowane
pakiety, pozwalające na porównanie wyników z literaturą:

\begin{itemize}
  \item \textbf{NIST SP 800-22} \cite{rukhin2010sp80022} - 15 testów
        obejmujących częstość, serie, rangę macierzy, transformatę Fouriera,
        złożoność liniową, entropię przybliżoną i inne; przeznaczony do
        oceny ciągów rzędu $10^6$ bitów i wyżej, z analizą rozkładu
        wartości $p$ dla wielu podciągów;
  \item \textbf{NIST SP 800-90B} \cite{turan2018sp80090b} - nie pakiet
        testów zgodności, lecz zestaw estymatorów min-entropii (ścieżka
        i.i.d. oraz ścieżka non-i.i.d., w której przyjmuje się minimum
        z dziesięciu estymatorów, m.in. kolizyjnego, kompresyjnego,
        markowskiego i opartych na predykcji);
  \item \textbf{Dieharder} oraz \textbf{TestU01} \cite{lecuyer2007testu01}
        --- rozbudowane baterie testów (w tym \emph{BigCrush}) o istotnie
        wyższej mocy, wymagające jednak strumieni rzędu $10^{11}$ bitów,
        co dla generatorów sprzętowych o umiarkowanej przepustowości bywa
        praktycznie nieosiągalne;
  \item \textbf{AIS-31} \cite{killmann2011ais31} - procedury T0-T8
        stosowane w certyfikacji europejskiej, obejmujące zarówno strumień
        surowy, jak i wyjściowy, wraz z wymaganiem istnienia
        stochastycznego modelu źródła.
\end{itemize}

Podział ról między analizą własną a pakietami standaryzowanymi jest zatem
następujący: testy dedykowane (obciążenie, autokorelacja, mapy 2D)
lokalizują \emph{przyczynę} defektu i pozwalają powiązać ją z parametrami
układu, natomiast pakiety standaryzowane orzekają o \emph{globalnej}
zgodności strumienia z modelem idealnym. Żadna z tych metod nie zastępuje
oszacowania min-entropii, która - zgodnie z porządkiem
\eqref{eq:porzadek} - pozostaje ostateczną miarą siły kryptograficznej
generatora.